\begin{pastexam}[label={px:ntu-110-comm-dsa-13}]{臺大電信所 110 學年度（B）第 13 題}
（單選；每題 2 分，答錯不倒扣）

Consider a divide-and-conquer algorithm, which solves a problem of size $n$ by dividing it into two subproblems of size $n/3$ and $n/5$, respectively. The solutions of the subproblems are then combined in $\Theta(n^2)$ time. Which of the following is correct about the time complexity $T(n)$ of this algorithm?
\begin{enumerate}[label=\Alph*.]
\item $T(n)$ must be in $\Theta(n^2)$.
\item $T(n)$ must be in $O(n^2)$ but not necessarily in $\Omega(n^2)$.
\item $T(n)$ must be in $\Omega(n^2)$ but not necessarily in $O(n^2)$.
\item None of the above. The complexity of $T(n)$ depends on the exact value of $n$.
\end{enumerate}
\end{pastexam}

\begin{pastsoln}
答案是 A。

遞迴式是 $T(n) = T(n/3) + T(n/5) + \Theta(n^2)$。 兩個子問題大小不同，主定理不能直接用，畫遞迴樹
（資結書第~2.3~節）。 設合併的工作是 $cn^2$：
\begin{itemize}
\item 第 $0$ 層：$cn^2$。
\item 第 $1$ 層：$c(n/3)^2 + c(n/5)^2 = cn^2\left(\tfrac19 + \tfrac1{25}\right) = \tfrac{34}{225}cn^2$。
\item 第 $k$ 層：每往下一層，每個節點的工作量都分成兩份，總和乘上 $\tfrac{34}{225}$，所以第 $k$ 層
至多 $\left(\tfrac{34}{225}\right)^k cn^2$（較淺的分支先到底，實際只會更少）。
\end{itemize}
各層加起來是公比 $\tfrac{34}{225} < 1$ 的等比級數：
\[
cn^2 \le T(n) \le cn^2 \sum_{k \ge 0} \left(\tfrac{34}{225}\right)^k = \tfrac{225}{191}\,cn^2.
\]
下界來自根節點本身就做了 $cn^2$。 所以 $T(n) = \Theta(n^2)$，選 A。 也可以用代換法證上界：
假設 $T(m) \le dm^2$ 對 $m < n$ 成立，$T(n) \le d\left(\tfrac{34}{225}\right)n^2 + cn^2 \le dn^2$，
只要 $d \ge \tfrac{225}{191}c$（資結書第~2.2~節）。

B 說不一定是 $\Omega(n^2)$，但根節點就花了 $\Theta(n^2)$，錯。 C 說不一定是 $O(n^2)$，上面的級數證明
了上界，錯。 D：$n/3$、$n/5$ 不是整數時取上下取整不影響漸近結果（資結書第~2.1.6~節），錯。
\end{pastsoln}
